% TOP_PROBLEM: {"id": 30006926, "title": "Colmezs Conjecture on Faltings Heights of CM Abelian Varieties", "category_id": 1, "category": {"id": 1, "name": "number_theory", "display_name": "Number Theory", "description": "Properties of integers, prime numbers, Diophantine equations.", "slug": "number-theory", "order_index": 1, "created_at": "2026-07-31T15:26:25.670Z"}, "status": "open"}
% FIRST_POSTED: 2026-09-25
% ENTRY_KIND: statement_only
% !TeX program = lualatex
% Definition and sources only; this file is not an LLM solution attempt.
\documentclass[11pt]{article}
\usepackage[margin=25mm]{geometry}
\usepackage{fontspec}
\setmainfont{Latin Modern Roman}
\usepackage{amsmath,amssymb,mathtools}
\usepackage{unicode-math}
\setmathfont{Latin Modern Math}
\usepackage{xurl,hyperref}
\hypersetup{colorlinks=true,urlcolor=blue}
\setcounter{secnumdepth}{0}
\setlength{\parindent}{0pt}
\setlength{\parskip}{5pt}
\setlength{\emergencystretch}{3em}
\begin{document}
\section*{\texorpdfstring{Colmezs Conjecture on Faltings Heights of CM Abelian Varieties}{Colmezs Conjecture on Faltings Heights of CM Abelian Varieties}}\label{30006926}
\subsection{Definitions and mathematical statement}
A CM field $E$ is a totally imaginary quadratic extension of a totally real field; write $[E:{\mathbb{Q}}]=2g$. A CM type $\Phi$ chooses one embedding from each complex-conjugate pair. For an abelian variety with full-ring complex multiplication of type $(\mathcal O_E,\Phi)$, use the stable Faltings height normalized as in Howard, Sections 3.2--3.3, \href{https://arxiv.org/html/1811.00428v2}{[E1]}: the normalized arithmetic degree of the line of invariant top differential forms, with the archimedean norm given there.

For $G_{{\mathbb{Q}}}=\operatorname{Gal}(\overline{{\mathbb{Q}}}/{\mathbb{Q}})$ define $A_\Phi(\sigma)=|\Phi\cap\sigma\Phi|$, and average this function over the finite $G_{{\mathbb{Q}}}$-orbit of $\Phi$, obtaining a class function $A_\Phi^0$. Expand it in irreducible Artin characters:
\[
 A_\Phi^0=\sum_\chi m_\Phi(\chi)\chi.
\]
With $L(s,\chi)$ excluding archimedean factors and $f_\chi$ its Artin conductor, the conjecture is
\[
 h^{\rm Falt}(E,\Phi)=
 -\sum_\chi m_\Phi(\chi)
 \left(\frac{L'(0,\chi)}{L(0,\chi)}+\frac12\log f_\chi\right).
\]
This is an identity for each CM type, not merely an average. The full-ring hypothesis and height normalization matter: heights for nonmaximal endomorphism orders require correction terms and are not silently included in this formula.
\par\smallskip\textit{Formulation and concept references:} \href{https://arxiv.org/abs/2603.28536}{[S1]}, \href{https://arxiv.org/html/1811.00428v2}{[E1]}.

\subsection{Short English statement}
For each complex multiplication type, does the Faltings height equal the precise expression built from Artin L-functions and their conductors?
\subsection{Sources}
\begin{itemize}
\item[S1]Maillot and Rossler, The Conjecture of Colmez and Reciprocity Laws for Modular Forms, abstract.
\url{https://arxiv.org/abs/2603.28536}.
\textit{Formulation source.}
\item[E1]Benjamin Howard, On the averaged Colmez conjecture, arXiv:1811.00428v2 (2018), Sections 3.2–3.3, Conjecture 3.3.1.
\url{https://arxiv.org/html/1811.00428v2}.
\textit{Added primary source: Definitions and normalization of the Faltings height, the CM-type class function, and the individual Colmez identity. Consulted 24 September 2026.}
\end{itemize}
\end{document}
